\scanpage{77}
% archive foliation: 76
\deleted{donc
\[
(\lambda^n-\lambda^{n'})g_{n'}=0\qquad\forall\lambda\in\mathbf Z^r,
\]
et il suffit de voir que les $\lambda^n-\lambda^{n'}$ engendrent l'idéal $1$, ou encore}

Si on exigeait qu'on a, pour tt $\mathbf Z$-algèbre $k$, des compositions entre les $\unsure{\mathcal P|P}_k(X,Y)=\unsure{\mathcal P|P}(X,Y)\otimes_{\mathbf Z}k$, et on exigeait la décomposition plus haut valable pour tt $\unsure{\mathcal P|P}_k(\prod X_i,Y)$, pour $\lambda\in(k^\times)^r$. Or on va mettre comme \uline{structure supplémentaire} la \uline{donnée} de la décomposition $(*)$ avec

e) Condition des degrés pour les compositions
\diagram{\begin{tikzcd}
X_1 \arrow[r,phantom,"\times\cdots\times"] & X_r \arrow[r,"{(n_1,\ldots,n_r)}"] & Y \\
Z_1^1\times\cdots\times Z_{s_1}^1 \arrow[u,"{m_1^1,\ldots,m_{s_1}^1}"] & Z_1^r,\ldots,Z_{s_r}^r \arrow[u,"{m_1^r,\ldots,m_{s_r}^r}"']
\end{tikzcd}}
la composée est de degré $(n_i m_1^i,n_i m_2^i,\ldots,n_i m_{s_i}^i)$ en $Z_1^i,\ldots,Z_{s_i}^i$; de même
\margin{Oubli: \unsure{Ces catégories|les catégories} \unsure{Ce qu'on veut est} \unsure{la structure} \unsure{de degré $1$} $\unsure{\unsure{\mathcal P|P}_1(X,Y)=\unsure{\mathcal P|P}_1(X,Y)}$.}
\margin{\deleted{\unsure{Pour tt $i$, $f_i=0$ aussi}} $\unsure{(n_i)_{i=1,\ldots,r}}$ \unsure{ou $n_i=0$}; \unsure{si $m_i=0$, $n_i\ne0$} \unsure{(\text{homogène})} \unsure{$\unsure{\mathcal P|P}_{(n_i)}$}.}

f) conditions \inserted{\unsure{symétrie}} pour la permutation des facteurs

g) $\forall X\in\operatorname{Ob}\unsure{\mathcal P|P}$, l'application canonique déduite de $X\longrightarrow0$
\[
\unsure{\mathcal P|P}_0(0,Y)\longrightarrow \unsure{\mathcal P|P}_0(X,Y)
\]
est bijective (les ``applications constantes'' \inserted{de $X$ en $Y$} ``ne dépendent pas de $X$'')

Si $k$ est un anneau comm. et si $\unsure{\mathcal P|P}_1$ est munie d'une structure $k$-linéaire par
\diagram{\begin{tikzcd}
k \arrow[r,no head] & \operatorname{End}(\unsure{\mathcal P|P}_1) & \operatorname{End}(\unsure{\mathcal P|P}) \arrow[l,no head]
\end{tikzcd}}
\deleted{on dit que} telle que les formules haut soient vraies pour $\lambda\in k^r$ (non seulmt $\lambda\in\mathbf Z^r$), on dit que $\unsure{\mathcal P|P}$ est $k$-polynomiale.

